%--------------------------------------------------------------------------
%		Não é necessário mudar qualquer coisa no preambulo deste arquivo
%		basta apenas colocar o seu texto, e as palavras-chave
%		observando as normas ABNT
%--------------------------------------------------------------------------
\chapter*{ABSTRACT} 
\thispagestyle{empty}

Irrigation accounts for about 70\% of the freshwater withdrawn worldwide, and scheduling based on fixed intervals or on the farmer's observation tends to apply more or less water than the crop requires. Commercial solutions that automate this decision are usually proprietary, expensive and offer little flexibility regarding the sensors that can be used. This work aimed to develop a low-cost smart irrigation system based on the Internet of Things, capable of monitoring the soil and the environment, estimating the crop water requirement and automating drip irrigation, with remote monitoring and control. This is applied and exploratory research, carried out incrementally, in which a first prototype based on a Raspberry Pi and MQTT gave way to the final architecture, organized into three layers. In the perception layer, a device with an ESP32, a soil moisture sensor, a DHT11 sensor, a relay module and an LCD display reads the sensors and decides when to irrigate, based on a moisture threshold, and how much to irrigate, using the FAO-56 method, with the reference evapotranspiration estimated by the Hargreaves-Samani equation. In the cloud layer, a REST API written in Java with Spring Boot receives the data, enforces the access rules, stores the information in Cloud Firestore and delivers commands to the device through long polling. In the application layer, an Angular web dashboard allows users to follow the readings and the history, configure the agronomic parameters and trigger irrigation, with distinct access profiles. The system was evaluated on a test bench through ten end-to-end test scenarios, supported by device logs, database documents, dashboard screens and API responses. The results showed that the readings reached the dashboard with the same values recorded by the device, that commands were executed in manual and automatic modes with the calculated durations, and that the safety locks and access control worked as expected. The tests also revealed adjustments to be made in the firmware, such as recording the volume of interrupted irrigations. It is concluded that it is feasible to build an automated and modular irrigation system with components estimated at between R\$~200 and R\$~345, although its agronomic efficiency still depends on a field evaluation.

%--------------------------------------------------------------------------
\vspace{2cm}
{\Large \textbf{Keywords:}}
\vspace{0.4cm}
%--------------------------------------------------------------------------

Internet of Things. Smart irrigation. ESP32. REST API. Evapotranspiration.
